\documentclass[10pt,a4paper]{amsart}
\usepackage[margin=19mm]{geometry}
\usepackage{amsmath,amssymb,mathtools}
\usepackage[scr=rsfs,cal=euler]{mathalfa}
\usepackage{xfrac}
\usepackage[unicode,hidelinks,bookmarks=false]{hyperref}
\newcommand{\ie}{i.e.\ }
\newcommand{\cf}{cf.\ }
\newcommand{\resp}{respectively\ }
\newcommand{\Subsection}{\subsection}
\providecommand{\nobreakdash}{\nobreak\hbox{-}}
\setlength{\emergencystretch}{2em}
\multlinegap0pt
\usepackage{braket}
\usepackage{multirow}
\usepackage{amssymb}
\usepackage{xcolor}
\usepackage{booktabs}
\usepackage{algorithm}\usepackage{algpseudocode}
\usepackage{tikz}
\usepackage{natbib}
\newtheorem{lemma}{Lemma}
\newcommand{\rV}{\mathsf{V}}
\title{A mathematical model for a universal digital quantum computer with an application to the Grover-Rudolph algorithm}
\author{Antonio Falc\'o, Daniela Falc\'o--Pomares, Hermann G. Matthies}
\date{Source: arXiv:2503.13388v5 (CC BY 4.0)}
\begin{document}
\maketitle

\noindent\emph{Source and licence.} A mathematical model for a universal digital quantum computer with an application to the Grover-Rudolph algorithm. Version arXiv:2503.13388v5 (CC BY 4.0), distributed by arXiv under the Creative Commons Attribution 4.0 International licence, http://creativecommons.org/licenses/by/4.0/, as recorded in the arXiv licence metadata for this exact version and shown on the abstract page https://arxiv.org/abs/2503.13388v5. Full licence notice: \texttt{licenses/arxiv-2503.13388-CC-BY-4.0.txt}.\par\smallskip

\noindent\textit{Editorial scope.} Counted unit: the article's lemma reduction (source0.tex lines 1944--1953). The article's complete original proof of it is delivered with it (source0.tex lines 1954--2059, one \texttt{proof} environment of 106 lines). No mathematics is written, repaired or re-derived: the statement and the proof are the article's own lines, in the article's own notation, and no retained line is substituted or re-flowed.  No further source-local context is needed: every cross-reference of every retained line resolves inside the two delivered spans.  Nothing was omitted from any retained span, so the package carries no omission record.  No retained line contains a citation, so the wrapper carries no bibliography and no bibliography entry.  No retained line contains a figure environment, an image-inclusion command or drawing code, so no image is delivered and the package declares no asset dependency.  Editorial formatting only, in addition: an AMS \texttt{amsart} preamble that restates the article's own declarations and macros used by the retained lines, namely the environments \texttt{lemma} on separate counters, together with the standard packages those lines need.  The retained environments print in this document as Lemma 1 in order of appearance, whereas the article prints the same items with the numbering of its own full document.  The complete native source of the article as distributed in the arXiv e-print for this exact version is bundled unchanged as \texttt{sources/174931/source0.tex}.
    \begin{lemma}\label{reduction}
      Let $\ket{\Psi} \in \mathbb{H}_N$ be a non-zero vector for some $N\ge 3.$ Then there exists 
      $$
      (\rV_1^{(N)},i_1^{(N)}),\ldots,(\rV_{N-1}^{N},i_{N-1}^{(N)}) \in \mathrm{U}(2)  \times \mathrm{Emb}(\mathrm{U}(2),\mathrm{U}(N))
      $$ 
      such that
      $$
      i_{N-1}^{(N)}(\rV_{N-1}^{(N)}) \cdots i_1^{(N)}(\rV_1^{(N)}) \ket{\Psi} =  \begin{bmatrix} \|\ket{\Psi}\| \\ 0 \\ \vdots  \\ 0 \end{bmatrix}
      $$
      \end{lemma}

      \begin{proof}
      We prove this lemma by induction on $N.$ First, we consider the case $N=3.$ Assume that $\ket{\Psi} = \begin{bmatrix} \Psi_1 \\ \Psi_2 \\ \Psi_3 \end{bmatrix}$ and suppose first that $\Psi_1 \neq 0.$ Then we take the matrix
      $$
      K_{12}^{(3)}\left( \begin{bmatrix}
        \frac{\overline{\Psi_1}}{\sqrt{ |\Psi_1|^2+|\Psi_2|^2}} & \frac{\overline{\Psi_2}}{\sqrt{ |\Psi_1|^2+|\Psi_2|^2}} \\
        \frac{\Psi_2}{\sqrt{ |\Psi_1|^2+|\Psi_2|^2}} & -\frac{\Psi_1}{\sqrt{ |\Psi_1|^2+|\Psi_2|^2}} \end{bmatrix} \right)  =\begin{bmatrix}
      \frac{\overline{\Psi_1}}{\sqrt{ |\Psi_1|^2+|\Psi_2|^2}} & \frac{\overline{\Psi_2}}{\sqrt{ |\Psi_1|^2+|\Psi_2|^2}} & 0 \\
      \frac{\Psi_2}{\sqrt{ |\Psi_1|^2+|\Psi_2|^2}} & -\frac{\Psi_1}{\sqrt{ |\Psi_1|^2+|\Psi_2|^2}} & 0 \\
      0 & 0 & 1
      \end{bmatrix}
      $$ 
      that satisfies
      $$
      \begin{bmatrix}
        \frac{\overline{\Psi_1}}{\sqrt{ |\Psi_1|^2+|\Psi_2|^2}} & \frac{\overline{\Psi_2}}{\sqrt{ |\Psi_1|^2+|\Psi_2|^2}} & 0 \\
        \frac{\Psi_2}{\sqrt{ |\Psi_1|^2+|\Psi_2|^2}} & -\frac{\Psi_1}{\sqrt{ |\Psi_1|^2+|\Psi_2|^2}} & 0 \\
        0 & 0 & 1
        \end{bmatrix} \begin{bmatrix} \Psi_1 \\ \Psi_2 \\ \Psi_3 \end{bmatrix} 
        = \begin{bmatrix} \sqrt{ |\Psi_1|^2+|\Psi_2|^2}\\ 0\\ \Psi_3\end{bmatrix}.
        .
      $$
      Now, we consider the matrix
      $$
      K_{13}^{(3)}\left( \begin{bmatrix}
        \frac{\sqrt{ |\Psi_1|^2+|\Psi_2|^2}}{\|\ket{\Psi}\|}  & \frac{\overline{\Psi_3}}{\|\ket{\Psi}\|} \\
        & \\
        \frac{\Psi_3}{\|\ket{\Psi}\|} & -\frac{\sqrt{ |\Psi_1|^2+|\Psi_2|^2}}{\|\ket{\Psi}\|}
        \end{bmatrix}\right) = \begin{bmatrix}
          \frac{\sqrt{ |\Psi_1|^2+|\Psi_2|^2}}{\|\ket{\Psi}\|} & 0  &\frac{\overline{\Psi_3}}{\|\ket{\Psi}\|} \\
        0 & 1 & 0 \\
        \frac{\Psi_3}{\|\ket{\Psi}\|} & 0 & -\frac{\sqrt{ |\Psi_1|^2+|\Psi_2|^2}}{\|\ket{\Psi}\|}
        \end{bmatrix},
      $$
      and hence 
      $$
      \begin{bmatrix}
        \frac{\sqrt{ |\Psi_1|^2+|\Psi_2|^2}}{\|\ket{\Psi}\|} & 0  &\frac{\overline{\Psi_3}}{\|\ket{\Psi}\|} \\
      0 & 1 & 0 \\
      \frac{\Psi_3}{\|\ket{\Psi}\|} & 0 & -\frac{\sqrt{ |\Psi_1|^2+|\Psi_2|^2}}{\|\ket{\Psi}\|}
      \end{bmatrix}\begin{bmatrix} \sqrt{ |\Psi_1|^2+|\Psi_2|^2}\\ 0\\ \Psi_3\end{bmatrix} = \begin{bmatrix} \|\ket{\Psi}\|\\ 0\\ 0\end{bmatrix}.
      $$
      Thus, there exists $(V_1^{(3)},i_1^{(3)}),(V_2^{(3)},i_2^{(3)}) \in \mathrm{U}(2) \times\mathrm{Emb}(\mathrm{U}(2),\mathrm{U}(3))$ such that 
      $$
      i_2^{(3)}(V_2^{(3)})i_1^{(3)}(V_1^{(3)})\ket{\Psi} = \begin{bmatrix} \|\ket{\Psi}\|\\ 0\\ 0\end{bmatrix}.
      $$
  The cases $\Psi_1 = 0,\Psi_2\neq 0$ and $\Psi_1 = 0,\Psi_2 =0,\Psi_3\neq 0$ are solved by using the product of matrices
  $$
  K_{12}^{(3)}\left( \begin{bmatrix} 0 & 1 \\ 1 & 0 \end{bmatrix}\right)K_{23}^{(3)}\left(\begin{bmatrix}
    \frac{\overline{\Psi_2}}{\sqrt{ |\Psi_2|^2+|\Psi_3|^2}} & \frac{\overline{\Psi_3}}{\sqrt{ |\Psi_2|^2+|\Psi_3|^2}} \\
    \frac{\Psi_3}{\sqrt{ |\Psi_2|^2+|\Psi_3|^2}} & -\frac{\Psi_2}{\sqrt{ |\Psi_2|^2+|\Psi_3|^2}} \end{bmatrix} \right),
  $$
  because
  \begin{align*}
  K_{12}^{(3)}\left( \begin{bmatrix} 0 & 1 \\ 1 & 0 \end{bmatrix}\right)K_{23}^{(3)}\left(\begin{bmatrix}
    \frac{\overline{\Psi_2}}{\sqrt{ |\Psi_2|^2+|\Psi_3|^2}} & \frac{\overline{\Psi_3}}{\sqrt{ |\Psi_2|^2+|\Psi_3|^2}} \\
    \frac{\Psi_3}{\sqrt{ |\Psi_2|^2+|\Psi_3|^2}} & -\frac{\Psi_2}{\sqrt{ |\Psi_2|^2+|\Psi_3|^2}} \end{bmatrix} \right)\begin{bmatrix} 0 \\ \Psi_2 \\ \Psi_3 \end{bmatrix} \\ 
      = K_{12}^{(3)}\left( \begin{bmatrix} 0 & 1 \\ 1 & 0 \end{bmatrix}\right) \begin{bmatrix} 0\\ \sqrt{ |\Psi_2|^2+|\Psi_3|^2}\\ 0 \end{bmatrix} = \begin{bmatrix} \|\ket{\Psi}\| \\ 0\\ 0 \end{bmatrix}.,
  \end{align*}
  
      In consequence the result follows for $N=3.$ Assume the result is true for $N-1$ and consider 
      $$
      \ket{\Psi} = \begin{bmatrix} \Psi_1 \\ \Psi_2 \\ \vdots \\ \Psi_{N} \end{bmatrix} = \begin{bmatrix}\ket{\Psi^{\prime}}\\ \Psi_N \end{bmatrix} \in \mathbb{H}_N, \text{ where } \ket{\Psi'} = \begin{bmatrix} \Psi_1 \\ \Psi_2 \\ \vdots \\ \Psi_{N-1} \end{bmatrix} \in \mathbb{H}_{N-1}
      $$ 
      be a non-zero vector. If $\ket{\Psi^{\prime}} \neq 0,$ we can apply the induction hypothesis to $\ket{\Psi^{\prime}}.$ This implies the existence of $(\rV_1^{(N-1)},i_1^{(N-1)}),\ldots,(\rV_{N-2}^{(N-1)},i_{N-2}^{(N-1)}) \in\mathrm{U}(2) \times\mathrm{Emb}(\mathrm{U}(2),\mathrm{U}(N-1)) $ such that 
      $$
      i_{N-2}^{(N-1)}(\rV_{N-2}^{(N-1)}) \cdots i_1^{(N-1)}(\rV_1^{(N-1)}) \ket{\Psi^{\prime}} =  \begin{bmatrix} \|\ket{\Psi^{\prime}}\| \\ 0 \\ \vdots  \\ 0 \end{bmatrix}.
      $$
    Hence, 
    \begin{align*}
    & \mathrm{up}_N(i_{N-2}^{(N-1)}(\rV_{N-2}^{(N-1)})) \cdots  \mathrm{up}_N(i_1^{(N-1)}(\rV_1^{(N-1)})) \ket{\Psi} \\ 
    & \\
    = & \mathrm{up}_N\left(i_{N-2}^{(N-1)}(\rV_{N-2}^{(N-1)}) \cdots  i_1^{(N-1)}(\rV_1^{(N-1)})\right)  \begin{bmatrix}\ket{\Psi^{\prime}}\\ \Psi_N \end{bmatrix}
    \\ 
    = & \begin{bmatrix} i_{N-2}^{(N-1)}(\rV_{N-2}^{(N-1)}) \cdots i_1^{(N-1)}(\rV_1^{(N-1)}) \ket{\Psi^{\prime}} \\ \Psi_N \end{bmatrix} = \begin{bmatrix} \|\ket{\Psi^{\prime}}\| \\ 0 \\ \vdots  \\ 0 \\ \Psi_N \end{bmatrix}.
    \end{align*}
    To conclude, we take the matrix 
    $$
  K_{1(N-1)}^{(N)}\left(
    \begin{bmatrix}
      \frac{\|\ket{\Psi^{\prime}}\|}{\|\ket{\Psi}\|} & \frac{\overline{\Psi_N}}{\|\ket{\Psi}\|} \\
      & \\
      \frac{\Psi_N}{\|\ket{\Psi}\|} & - \frac{\|\ket{\Psi^{\prime}}\|}{\|\ket{\Psi}\|}
    \end{bmatrix}
  \right)\begin{bmatrix} \|\ket{\Psi^{\prime}}\| \\ 0 \\ \vdots  \\ 0 \\ \Psi_N \end{bmatrix} = \begin{bmatrix} \|\ket{\Psi}\| \\ 0 \\ \vdots  \\ 0 \\ 0 \end{bmatrix},
    $$
  here we use that $\|\ket{\Psi}\|^2 =\|\ket{\Psi^{\prime}}\|^2 + |\Psi_N|^2.$ In consequence we have that 
  $$
  K_{1(N-1)}^{(N)}(V_{N-1}^{(N-1)})(\mathrm{up}_N \circ i_{N-2}^{(N-1)})(\rV_{N-2}^{(N-1)})) \cdots  (\mathrm{up}_N \circ i_1^{(N-1)})(\rV_1^{(N-1)})) \ket{\Psi} = \begin{bmatrix} \|\ket{\Psi}\| \\ 0 \\ \vdots  \\ 0 \\ 0 \end{bmatrix},
  $$
  where 
  $
  (V_{N-1}^{(N-1)},K_{1(N-1)}^{(N)}),(\rV_{N-2}^{(N-1)},\mathrm{up}_N \circ i_{N-2}^{(N-1)}),\ldots,(\rV_1^{(N-1)},\mathrm{up}_N \circ i_1^{(N-1)})$ are in  $\mathrm{U}(2) \times\mathrm{Emb}(\mathrm{U}(2),\mathrm{U}(N)) .$ Otherwise, if $\ket{\Psi^{\prime}} = 0,$ then $\Psi_N \neq 0,$ and we have
  \begin{align*}
    K_{12}^{(N)}\left( \begin{bmatrix}
      0 & 1 \\
     1 & 0 \end{bmatrix}
     \right) \cdots K_{N-2(N-1)}^{(N)}\left( \begin{bmatrix}
      0 & 1 \\
     1 & 0 \end{bmatrix}
     \right) K_{N(N-1)}^{(N)}\left( \begin{bmatrix}
     0 & \frac{\overline{\Psi_N}}{|\Psi_N|} \\
    \frac{\Psi_N}{|\Psi_N|} & 0 \end{bmatrix}
    \right)\begin{bmatrix} 0 \\ 0 \\ \vdots \\ 0 \\ \Psi_{N} \end{bmatrix} = \begin{bmatrix}  |\Psi_N| \\ 0 \\ \vdots \\ 0 \\ 0 \end{bmatrix}.
    \end{align*}
  This concludes the proof of the lemma.
  \end{proof}

\end{document}
